\documentclass[aspectratio=169]{beamer}
\usetheme{Madrid}
\usecolortheme{default}
\setbeamertemplate{footline}[frame number]

\usepackage[utf8]{inputenc}
\usepackage[T1]{fontenc}
\usepackage{graphicx}
\usepackage{booktabs}
\usepackage{amsmath}
\usepackage{amssymb}
\usepackage{array}

\title[190 MeV/u d on Sn124]{190 MeV/u Polarized Deuteron on $^{124}$Sn\\Full-Chain Rate Estimate}
\subtitle{ImQMD shell-sum cross sections and Geant4 coincidence rates}
\author{Tian}
\date{2026-04-12}

\begin{document}

\begin{frame}
\titlepage
\end{frame}

\begin{frame}{Dataset and Analysis Scope}
\small
\begin{itemize}
  \item Reaction considered here: polarized deuteron on $^{124}$Sn at 190 MeV/u.
  \item The quantitative results are based on the current analysis chain:
  \begin{itemize}
    \item ImQMD raw breakup sample in \texttt{z\_pol/b\_discrete},
    \item fresh Geant4 rerun in the current detector geometry,
    \item rate conversion using the recomputed $^{124}$Sn number density.
  \end{itemize}
  \item Event sample used for the numbers shown below:
  \begin{itemize}
    \item target: $^{124}$Sn,
    \item beam energy: 190 MeV/u,
    \item sample: z polarization, $\gamma = 0.60$,
    \item impact-parameter range used in the shell sum: $b = 1 \ldots 9$ fm.
  \end{itemize}
\end{itemize}
\end{frame}

\begin{frame}{Beam and Target Inputs}
\small
\begin{itemize}
  \item Beam intensity: $I = 1.0\times 10^7$ pps
  \item Beam time for the headline yield: $t = 16$ h
  \item Beam energy per nucleon: $T_u = 190$ MeV/u
  \item SRC harmonic number: $h = 6$
  \item SRC extraction radius: $R_{\mathrm{ext}} = 5.36$ m
  \item Tin density used for target conversion: $\rho = 7.31~\mathrm{g/cm^3}$
  \item Enriched-$^{124}$Sn number density:
  \[
  n = \frac{N_A \rho}{M_{124}} =
  \frac{(6.022\times 10^{23})(7.31)}{124}
  = 3.5501\times 10^{22}~\mathrm{cm^{-3}}
  \]
\end{itemize}

\vspace{0.1cm}
\small
The full-chain cross sections extracted from \texttt{validation\_summary.json} are:
\begin{itemize}
  \item selected-input cross section: $\sigma_{\mathrm{sel}} = 516.5$ mb
  \item detected-coincidence cross section: $\sigma_{\mathrm{coin}} = 204.2$ mb
\end{itemize}
\end{frame}

\begin{frame}{SRC Bunch Timing at 190 MeV/u}
\small
\begin{align*}
\gamma &= 1 + \frac{T_u}{m_u c^2}
       = 1 + \frac{190}{931.494}
       = 1.20397, \\
\beta &= \sqrt{1-\gamma^{-2}}
      = 0.55689, \\
v &= \beta c
  = 1.6695\times 10^8~\mathrm{m/s}, \\
C &= 2\pi R_{\mathrm{ext}}
  = 2\pi \times 5.36
  = 33.678~\mathrm{m}, \\
f_{\mathrm{rev}} &= \frac{v}{C}
                 = 4.9573~\mathrm{MHz}, \\
f_{\mathrm{bucket}} &= h f_{\mathrm{rev}}
                    = 29.7441~\mathrm{MHz}, \\
\Delta t &= \frac{1}{f_{\mathrm{bucket}}}
         = 33.62~\mathrm{ns}.
\end{align*}

\[
\mu_{\mathrm{beam/bucket}} = \frac{I}{f_{\mathrm{bucket}}} \approx 0.3362
\]

\small
The bunch interval is therefore fixed by the beam kinematics and the SRC machine parameters.
\end{frame}

\begin{frame}{Cross-Section Extraction from ImQMD and Geant4}
\tiny
\textbf{Step 1: breakup probability from ImQMD in each impact-parameter shell}
\[
P_{\mathrm{bu}}(b_i)=\frac{N_{\mathrm{breakup}}(b_i)}{N_{\mathrm{incident}}(b_i)}
\]

\textbf{Step 2: shell-sum cross section for the current \texttt{b\_discrete} dataset}
\[
\sigma \approx \sum_i 2\pi b_i \Delta b\, P(b_i)
\]

\textbf{Step 3: propagate the selected events through Geant4}
\[
P_{\mathrm{sel}}(b_i)=\frac{N_{\mathrm{G4}}(b_i)}{N_{\mathrm{incident}}(b_i)},\qquad
P_{\mathrm{coin}}(b_i)=\frac{N_{\mathrm{coin}}(b_i)}{N_{\mathrm{incident}}(b_i)}
\]

\[
\sigma_{\mathrm{sel}} \approx \sum_i 2\pi b_i \Delta b\, P_{\mathrm{sel}}(b_i),\qquad
\sigma_{\mathrm{coin}} \approx \sum_i 2\pi b_i \Delta b\, P_{\mathrm{coin}}(b_i)
\]

\vspace{0.1cm}
\textbf{Coincidence definition}
\begin{itemize}
  \item The detector geometry is defined by the current Geant4 macro.
  \item The coincidence rate quoted here comes from \textbf{Geant4 event counting} in that geometry, rather than from a geometry-only solid-angle estimate.
  \item In the current validation script, a coincidence event means the same event has both:
  \begin{itemize}
    \item at least one PDC fragment hit,
    \item at least one NEBULA hit.
  \end{itemize}
\end{itemize}
\end{frame}

\begin{frame}{Full-Chain Event Accounting}
\small
\centering
\begin{tabular}{>{\raggedright\arraybackslash}p{4.4cm}cc}
\toprule
Quantity & Value & Comment \\
\midrule
Raw-file header incident events & 180000 & 18 files, 10000 events each \\
Stored breakup rows before cut & 28611 & breakup events present in raw files \\
Events propagated through Geant4 & 22174 & after generator selection and cut \\
Fresh coincidence hits & 8408 & Geant4 rerun result \\
Fresh coincidence efficiency on propagated sample & 37.92\% & $8408 / 22174$ \\
Archive coincidence efficiency & 37.95\% & archived rerun for comparison \\
Fresh minus archive & $-0.027$ points & reproducibility check \\
\bottomrule
\end{tabular}

\vspace{0.2cm}
\small
The Geant4 coincidence efficiency is stable at the level of about $0.03$ percentage points between the fresh rerun and the archived output.
\end{frame}

\begin{frame}{Cross Sections from the Current Chain}
\small
\centering
\begin{tabular}{>{\raggedright\arraybackslash}p{4.8cm}cc}
\toprule
Quantity & Value & Meaning \\
\midrule
Raw breakup cross section & 678.2 mb & ImQMD breakup shell sum \\
Selected-input cross section & 516.5 mb & events that enter the current Geant4 chain \\
Detected-coincidence cross section & 204.2 mb & events counted as PDC+NEBULA coincidence \\
Naive $\pi b_{\max}^2 (N/N_0)$ estimate & 404.5 mb & shown only as a comparison \\
\bottomrule
\end{tabular}

\vspace{0.25cm}
\begin{itemize}
  \item For the present \texttt{b\_discrete} data structure, the shell-sum estimate is the appropriate quantity for rate conversion.
  \item The rate conversion is
  \[
  R_{\mathrm{sel}} = n d I \sigma_{\mathrm{sel}},\qquad
  R_{\mathrm{coin}} = n d I \sigma_{\mathrm{coin}}
  \]
  \item Here $\sigma_{\mathrm{coin}}$ already includes the current Geant4 detector acceptance and coincidence definition.
\end{itemize}
\end{frame}

\begin{frame}{Thickness Scan: Local Full-Chain Results for $^{124}$Sn}
\small
\centering
\begin{tabular}{>{\raggedright\arraybackslash}p{1.55cm}cccc}
\toprule
Thickness & $R_{\mathrm{sel}}$ & $R_{\mathrm{coin}}$ & 16 h coincidence yield & $\mu_{\mathrm{coin/bucket}}$ \\
(mm) & (kHz) & (kHz) &  &  \\
\midrule
0.5 & 9.169 & 3.625 & $2.09\times 10^8$ & $1.22\times 10^{-4}$ \\
1.0 & 18.338 & 7.250 & $4.18\times 10^8$ & $2.44\times 10^{-4}$ \\
1.5 & 27.506 & 10.875 & $6.26\times 10^8$ & $3.66\times 10^{-4}$ \\
2.0 & 36.675 & 14.500 & $8.35\times 10^8$ & $4.87\times 10^{-4}$ \\
3.0 & 55.013 & 21.750 & $1.25\times 10^9$ & $7.31\times 10^{-4}$ \\
\bottomrule
\end{tabular}

\vspace{0.3cm}
\begin{itemize}
  \item Even with the local full-chain coincidence cross section, the single-bucket coincidence occupancy remains well below $10^{-3}$.
  \item Under the current local chain, the 3 mm reference gives:
  \begin{itemize}
    \item $R_{\mathrm{coin}} \approx 21.75$ kHz
    \item $N_{\mathrm{coin}}(16~\mathrm{h}) \approx 1.25\times 10^9$
  \end{itemize}
\end{itemize}
\end{frame}

\begin{frame}{Inventory-Constrained Disk Designs}
\small
\centering
\begin{tabular}{>{\raggedright\arraybackslash}p{1.8cm}cccc}
\toprule
Disk diameter & Areal density & Bulk thickness & $R_{\mathrm{coin}}$ & 16 h coincidence yield \\
(mm) & (mg/cm$^2$) & (mm) & (kHz) &  \\
\midrule
12 & 1061 & 1.451 & 10.523 & $6.06\times 10^8$ \\
15 & 679 & 0.929 & 6.735 & $3.88\times 10^8$ \\
20 & 382 & 0.523 & 3.788 & $2.18\times 10^8$ \\
\bottomrule
\end{tabular}

\vspace{0.3cm}
\begin{itemize}
  \item All three options assume the same inventory constraint: 1.2 g of tin per isotope.
  \item The 15 mm disk remains a practical compromise:
  \begin{itemize}
    \item wider alignment tolerance than 12 mm,
    \item still much higher rate than 20 mm.
  \end{itemize}
\end{itemize}
\end{frame}

\begin{frame}{Interpretation of the Coincidence Rate}
\small
\begin{itemize}
  \item The coincidence rate presented here is a \textbf{full-chain} quantity:
  \[
  \mathrm{ImQMD} \rightarrow \mathrm{event~selection} \rightarrow \mathrm{Geant4} \rightarrow \mathrm{PDC+NEBULA~coincidence~counting}
  \]
  \item Its ingredients are:
  \begin{itemize}
    \item the breakup production part comes from ImQMD shell-by-shell probabilities,
    \item the detector part comes from Geant4 transport through the current detector geometry,
    \item the final coincidence rate uses the detected-coincidence cross section $\sigma_{\mathrm{coin}} = 204.2$ mb.
  \end{itemize}
  \item In other words, the quoted rate already combines the production model and the detector response within one consistent chain.
\end{itemize}
\end{frame}

\begin{frame}{Summary}
\small
\begin{itemize}
  \item Bunch spacing at 190 MeV/u in the current SRC setting:
  \[
  \Delta t = 33.62~\mathrm{ns}
  \]
  \item For a 3 mm $^{124}$Sn reference target:
  \[
  R_{\mathrm{coin}} \approx 21.75~\mathrm{kHz},\qquad
  N_{\mathrm{coin}}(16~\mathrm{h}) \approx 1.25\times 10^9
  \]
  \item For the 15 mm diameter, 1.2 g inventory-constrained disk:
  \[
  d \approx 0.929~\mathrm{mm},\qquad
  R_{\mathrm{coin}} \approx 6.735~\mathrm{kHz},\qquad
  N_{\mathrm{coin}}(16~\mathrm{h}) \approx 3.88\times 10^8
  \]
  \item Under the present detector definition, the coincidence rate is a Geant4-counted full-chain result derived from the current $^{124}$Sn zpol $\gamma=0.60$ sample.
\end{itemize}
\end{frame}

\end{document}
